\documentclass[a4paper,12pt]{article}

\usepackage[T2A]{fontenc}
\usepackage[utf8]{inputenc}
\usepackage[russian]{babel}
\usepackage{amsmath, amssymb, amsthm}
\usepackage{geometry}
\usepackage{tikz}
\usepackage{booktabs}
\usepackage{hyperref}
\usepackage{algorithm}
\usepackage{algorithmicx}
\usepackage{algpseudocode}

\geometry{top=2cm, bottom=2cm, left=2.5cm, right=1.5cm}

\newtheorem{theorem}{Теорема}
\newtheorem{lemma}{Лемма}
\newtheorem{definition}{Определение}

\title{Билет №77 \\ \small Основные принципы управления транзакциями, журналированием и восстановлением. (СпБГУ)}
\author{Сериков Владислав Александрович}
\date{16 мая 2026}

\begin{document}

\maketitle
\tableofcontents
\newpage

\section{Введение в теорию транзакций}

\begin{definition}
\textbf{Транзакция} — это логическая единица работы с базой данных (БД), состоящая из последовательности операций над данными (чтение, запись), которая должна быть выполнена целиком и успешно, либо не выполнена вообще.
\end{definition}

Для обеспечения надежности и консистентности данных транзакции должны удовлетворять свойствам \textbf{ACID}:
\begin{itemize}
    \item \textbf{A}tomicity (Атомарность): Транзакция выполняется полностью или не выполняется вовсе.
    \item \textbf{C}onsistency (Согласованность): Транзакция переводит БД из одного согласованного состояния в другое.
    \item \textbf{I}solation (Изолированность): Параллельно выполняемые транзакции не должны влиять друг на друга так, как если бы они выполнялись последовательно.
    \item \textbf{D}urability (Долговечность): Результаты успешно зафиксированной транзакции (COMMITTED) сохраняются в БД даже в случае аппаратных сбоев.
\end{itemize}

Управление транзакциями разбивается на две крупные подзадачи: управление параллельным выполнением (Concurrency Control) для обеспечения свойства Изоляции, и управление восстановлением (Recovery Management) для обеспечения свойств Атомарности и Долговечности.

\section{Управление параллельным выполнением}

\begin{definition}
\textbf{История (Расписание, Schedule)} $S$ — это хронологическая последовательность операций нескольких транзакций $T_1, \dots, T_n$.
\end{definition}

\begin{definition}
Две операции называются \textbf{конфликтующими}, если они принадлежат разным транзакциям, обращаются к одному и тому же элементу данных, и по крайней мере одна из них является операцией записи ($W$).
Типы конфликтов: Read-Write (RW), Write-Read (WR), Write-Write (WW).
\end{definition}

\begin{definition}
\textbf{Сериализуемое расписание} — это расписание, результат выполнения которого эквивалентен результату некоторого последовательного (serial) выполнения тех же транзакций.
Если эквивалентность достигается путем перестановки неконфликтующих операций, расписание называется \textbf{конфликтно-сериализуемым} (Conflict Serializable).
\end{definition}

\subsection{Граф предшествования (Serialization Graph)}

Для определения конфликтной сериализуемости строится ориентированный граф предшествования $G(S) = (V, E)$, где:
\begin{itemize}
    \item $V$ — множество зафиксированных транзакций в $S$.
    \item $E$ — направленные ребра. Ребро $T_i \to T_j$ существует, если в $S$ есть операция $p \in T_i$, конфликтующая с операцией $q \in T_j$, и $p$ выполняется раньше $q$.
\end{itemize}

\begin{theorem}[О конфликтной сериализуемости]
\label{thm:serial}
Расписание $S$ является конфликтно-сериализуемым тогда и только тогда, когда его граф предшествования $G(S)$ ацикличен.
\end{theorem}

\begin{proof}
\textbf{Необходимость ($\Rightarrow$):} Пусть $S$ конфликтно-сериализуемо. По определению существует эквивалентное последовательное расписание $S'$. Если в графе $G(S)$ есть ребро $T_i \to T_j$, это означает, что $T_i$ выполнила конфликтующую операцию раньше $T_j$. Следовательно, в любом эквивалентном последовательном расписании $S'$, транзакция $T_i$ должна предшествовать $T_j$. Предположим, что в $G(S)$ существует цикл $T_1 \to T_2 \to \dots \to T_k \to T_1$. Это означало бы, что в последовательном расписании $S'$ транзакция $T_1$ должна выполниться раньше $T_1$, что невозможно. Значит, граф ацикличен.

\textbf{Достаточность ($\Leftarrow$):} Пусть $G(S)$ ацикличен. Поскольку это направленный ациклический граф (DAG), его вершины можно топологически отсортировать. Пусть $T_{s_1}, T_{s_2}, \dots, T_{s_n}$ — топологическая сортировка транзакций. Сформируем последовательное расписание $S'$, в котором транзакции выполняются ровно в этом порядке. В расписании $S$ можно менять местами соседние неконфликтующие операции разных транзакций, не меняя итогового результата. Поскольку порядок всех конфликтующих операций совпадает с топологической сортировкой графа $G(S)$, мы можем с помощью таких перестановок привести $S$ к $S'$. Значит $S$ эквивалентно последовательному расписанию $S'$, то есть является конфликтно-сериализуемым.
\end{proof}

\subsection{Двухфазная блокировка (Two-Phase Locking, 2PL)}

Для автоматического обеспечения сериализуемости СУБД используют алгоритмы контроля конкурентного доступа, самым популярным из которых является 2PL.

\textbf{Правила протокола 2PL:}
\begin{enumerate}
    \item Перед чтением или записью объекта транзакция должна получить соответствующую блокировку (Shared для чтения, Exclusive для записи).
    \item Транзакция не может получить новую блокировку после того, как она освободила хотя бы одну блокировку.
\end{enumerate}

Таким образом, транзакция имеет фазу роста (только захват блокировок) и фазу сжатия (только снятие блокировок).

\begin{theorem}[Корректность 2PL]
Любое расписание $S$, сгенерированное протоколом 2PL, является конфликтно-сериализуемым.
\end{theorem}

\begin{proof}
Предположим от противного, что расписание $S$, порожденное 2PL, не является конфликтно-сериализуемым. Согласно Теореме \ref{thm:serial}, граф $G(S)$ содержит цикл: $T_1 \to T_2 \to \dots \to T_k \to T_1$.

Ребро $T_i \to T_{i+1}$ означает наличие конфликта, где $T_i$ обращалась к элементу $X$ раньше $T_{i+1}$. Следовательно, $T_i$ владела блокировкой на $X$ и должна была освободить ее (Unlock) до того, как $T_{i+1}$ смогла её захватить (Lock). Обозначим момент времени получения первой блокировки как $L_{first}(T)$, а снятия первой блокировки (точку фиксации блокировок, lock point) как $U_{first}(T)$.
Так как $T_i$ освободила блокировку на $X$ до того, как $T_{i+1}$ ее захватила, то:
$$ U_{first}(T_i) < L_{last}(T_{i+1}) $$
По правилу 2PL транзакция не может захватывать блокировки после освобождения любой блокировки, следовательно, $L_{last}(T_{i+1}) \le U_{first}(T_{i+1})$. Объединяя неравенства, получаем:
$$ U_{first}(T_i) < U_{first}(T_{i+1}) $$
Пройдя по всему циклу, мы получим:
$$ U_{first}(T_1) < U_{first}(T_2) < \dots < U_{first}(T_k) < U_{first}(T_1) $$
Следовательно, $U_{first}(T_1) < U_{first}(T_1)$, что является очевидным противоречием. Значит, предположение о наличии цикла неверно, и $S$ конфликтно-сериализуемо.
\end{proof}


\section{Журналирование (Logging) и упреждающая запись (WAL)}

Для восстановления после сбоев используется журнал транзакций (Log).
Основной принцип журналирования — \textbf{Write-Ahead Logging (WAL)}:
\begin{enumerate}
    \item Любое изменение данных сначала фиксируется в журнале. Запись в журнале должна быть сохранена на жесткий диск (flush) \textit{до} того, как измененная страница данных будет записана на диск.
    \item Все журнальные записи, относящиеся к транзакции, должны быть принудительно сохранены на диск \textit{до} того, как транзакция будет считаться успешно зафиксированной (COMMITTED).
\end{enumerate}

Каждая запись в журнале имеет уникальный номер \textbf{LSN (Log Sequence Number)}, который монотонно возрастает.
Типовая структура журнальной записи об обновлении содержит:
$[LSN, prevLSN, transID, type, pageID, undoNextLSN, old\_data, new\_data]$.

СУБД классифицируются по политикам управления буфером (Buffer Management):
\begin{itemize}
    \item \textbf{Steal:} Разрешается записывать на диск измененные страницы незафиксированных транзакций. (Требует хранения информации для UNDO).
    \item \textbf{No-Force:} При коммите транзакции не требуется немедленно сбрасывать все измененные ею страницы данных на диск. (Требует хранения информации для REDO).
\end{itemize}
Большинство современных СУБД используют комбинацию \textbf{Steal/No-Force}, так как она обеспечивает наивысшую производительность.

\section{Алгоритм восстановления ARIES}

\textbf{ARIES (Algorithm for Recovery and Isolation Exploiting Semantics)} — классический алгоритм восстановления, основанный на политике Steal/No-Force и протоколе WAL.

\subsection{Шаги алгоритма ARIES}
Процесс восстановления после сбоя (Crash Recovery) проходит в три фазы:

\begin{enumerate}
    \item \textbf{Фаза Анализа (Analysis Phase):}
    \begin{itemize}
        \item \textit{Цель:} Восстановить состояние СУБД на момент сбоя (список активных транзакций и список грязных страниц).
        \item \textit{Шаги:} Алгоритм читает журнал вперед, начиная с последнего успешного чекпоинта (Checkpoint).
        \item Если встречается запись \texttt{UPDATE} от транзакции $T_i$, $T_i$ добавляется в список транзакций \texttt{Transaction Table (TT)}, а измененная страница $P_j$ — в список грязных страниц \texttt{Dirty Page Table (DPT)} (с фиксацией $recLSN$ — номера записи, впервые загрязнившей страницу).
        \item Если встречается \texttt{COMMIT} или \texttt{ABORT}, статус в \texttt{TT} обновляется. Записи \texttt{END} удаляют транзакцию из \texttt{TT}.
    \end{itemize}

    \item \textbf{Фаза Повторения (Redo Phase):}
    \begin{itemize}
        \item \textit{Цель:} Воспроизвести историю базы данных (Repeating History), приведя данные к состоянию на момент сбоя, включая операции транзакций, которые потом будут отменены.
        \item \textit{Шаги:} Алгоритм определяет минимальный $recLSN$ среди всех страниц в DPT. С этого LSN журнал читается вперед до конца.
        \item Для каждой записи типа \texttt{UPDATE} алгоритм проверяет, нужно ли применять изменения. Изменение \textbf{не} применяется, если:
        1) Страницы нет в DPT;
        2) Страница есть в DPT, но ее $recLSN > LSN$ текущей записи;
        3) $pageLSN$ страницы на диске $\ge LSN$ текущей записи.
        \item Если условия не выполняются, страница загружается в память, к ней применяется \texttt{new\_data}, и ее $pageLSN$ обновляется до $LSN$ записи.
    \end{itemize}

    \item \textbf{Фаза Отмены (Undo Phase):}
    \begin{itemize}
        \item \textit{Цель:} Отменить действия всех транзакций, которые были активны на момент сбоя (находятся в TT) — так называемые транзакции-лузеры.
        \item \textit{Шаги:} Алгоритм просматривает журнал в обратном порядке (с конца). Для каждой транзакции из TT берется ее последний $LSN$.
        \item Выбирается запись с наибольшим $LSN$ среди всех подлежащих отмене.
        \item При отмене операции \texttt{UPDATE} (используя \texttt{old\_data}) в журнал пишется \textbf{CLR (Compensation Log Record)} — компенсационная запись. CLR содержит поле $undoNextLSN$, указывающее на предыдущую запись отменяемой транзакции ($prevLSN$). Это гарантирует, что в случае повторного сбоя во время восстановления, отмененные операции не будут отменяться дважды.
        \item Процесс повторяется, пока для всех транзакций-лузеров не будет достигнуто начало ($prevLSN = NULL$). Затем для каждой из них пишется запись \texttt{END}.
    \end{itemize}
\end{enumerate}

\subsection{Минимальный пример восстановления ARIES}

\textbf{Дано:} Исходно страницы $P1$ и $P2$ синхронизированы на диске.

\begin{table}[h]
\centering
\begin{tabular}{cll}
\toprule
\textbf{LSN} & \textbf{Запись в журнале} & \textbf{Комментарий} \\
\midrule
10 & \texttt{UPDATE T1, P1 (old: A, new: B)} & $T_1$ меняет $P1$ \\
20 & \texttt{UPDATE T2, P2 (old: X, new: Y)} & $T_2$ меняет $P2$ \\
30 & \texttt{COMMIT T1} & $T_1$ фиксируется. Журнал до LSN=30 пишется на диск \\
\bottomrule
\multicolumn{3}{c}{\textbf{CRASH (Сбой системы)}} \\
\end{tabular}
\end{table}

Предположим, из-за No-Force измененная $P1$ не была сброшена на диск. Из-за Steal $P2$ могла быть, а могла и не быть сброшена (предположим, что не сброшена). Чекпоинтов не было (читаем с $LSN=10$).

\textbf{Иллюстрация шагов восстановления:}

1. \textbf{Analysis Phase:}
\begin{itemize}
    \item Читаем LSN 10: \texttt{TT} = \{$T_1$\}, \texttt{DPT} = \{($P1, recLSN=10$)\}
    \item Читаем LSN 20: \texttt{TT} = \{$T_1, T_2$\}, \texttt{DPT} = \{($P1, recLSN=10$), ($P2, recLSN=20$)\}
    \item Читаем LSN 30: $T_1$ зафиксирована. \texttt{TT} = \{$T_2$\}.
    \item Итог: $T_2$ — активна (требует UNDO). DPT содержит $P1, P2$. Мин. $recLSN = 10$.
\end{itemize}

2. \textbf{Redo Phase:}
\begin{itemize}
    \item Старт с LSN 10.
    \item LSN 10: $P1$ на диске имеет $pageLSN < 10$. Применяем \texttt{new: B}.
    \item LSN 20: $P2$ на диске имеет $pageLSN < 20$. Применяем \texttt{new: Y}.
    \item LSN 30: Игнорируется (не операция с данными).
    \item База данных приведена в состояние на момент сбоя.
\end{itemize}

3. \textbf{Undo Phase:}
\begin{itemize}
    \item В \texttt{TT} только $T_2$. Последний LSN для $T_2$ — 20.
    \item Читаем LSN 20: Отменяем изменение (пишем в $P2$ значение \texttt{old: X}).
    \item Создаем компенсационную запись: \texttt{CLR: UNDO T2 LSN 20 (undoNextLSN=NULL)}.
    \item Так как больше записей у $T_2$ нет ($prevLSN$ для LSN 20 был NULL), транзакция полностью отменена. В журнал пишется \texttt{END T2}.
\end{itemize}

В результате $T_1$ зафиксирована (ее изменения в $P1$ восстановлены), а незафиксированная $T_2$ отменена. Консистентность базы данных полностью соблюдена.

\end{document}
